\section{Билет 15. Метод штрафных функций}

\ticketq{Метод штрафных функций.}

\paragraph{Постановка.}
\[
  \varphi(x) \to \min , \qquad
  x \in S = \{\, x \mid \psi_i(x) \le 0,\ i = \overline{1,m};\
    \chi_j(x) = 0,\ j = \overline{1,p} \,\} .
\]
Метод \textbf{внешних} штрафов сводит задачу к последовательности задач
\emph{безусловной} минимизации: нарушение ограничений не запрещается, а
штрафуется. Точки при этом могут быть недопустимыми.

\paragraph{Штрафная функция.}
$\alpha(x) \ge 0$, причём $\alpha(x) = 0 \iff x \in S$. Стандартные формы
($p > 1$, обычно $p = 2$):
\[
  \text{одно неравенство } \psi \le 0 :\quad
  \alpha(x) = \begin{cases} 0, & \psi(x) \le 0 , \\
                            \psi(x)^p, & \psi(x) > 0 ; \end{cases}
\]
\[
  \text{одно равенство } \chi = 0 :\quad \alpha(x) = \abs{\chi(x)}^p ;
\]
\[
  \text{общий случай:}\quad
  \alpha(x) = \sum_{i=1}^{m}\xi_i \max\{0, \psi_i(x)\}^p
            + \sum_{j=1}^{p}\eta_j \abs{\chi_j(x)}^p ,
  \qquad \xi_i, \eta_j > 0 .
\]

\paragraph{Вспомогательная задача.}
\[
  \varphi(x) + \mu\,\alpha(x) \to \min_{x \in V} , \quad \mu > 0 ,
  \qquad x_\mu \ \text{— решение} ,
  \qquad P(\mu) = \varphi(x_\mu) + \mu\,\alpha(x_\mu) .
\]

\paragraph{Лемма (свойства вспомогательных задач).}
\textit{Пусть все функции непрерывны и при каждом $\mu > 0$ решение $x_\mu$
существует. Тогда}
\begin{csteps}
  \item $\min\limits_{x\in S}\varphi(x) \ \ge\ \sup\limits_{\mu>0}P(\mu)$;
  \item $\alpha(x_\mu)$ — невозрастающая функция $\mu$;
  \item $\varphi(x_\mu)$ — неубывающая функция $\mu$;
  \item $P(\mu)$ — неубывающая функция $\mu$;
  \item $\lim\limits_{\mu\to\infty}P(\mu) = \sup\limits_{\mu>0}P(\mu)$.
\end{csteps}

\textsc{Док-во.}
\textbf{1)} Для $x \in S$ штраф нулевой, а $x_\mu$ минимизирует
$\varphi + \mu\alpha$ на $V \supset S$:
\[
  P(\mu) = \min_{z\in V}\big(\varphi(z)+\mu\alpha(z)\big)
   \le \varphi(x) + \mu\underbrace{\alpha(x)}_{=\,0} = \varphi(x) ;
\]
минимум по $x \in S$ и супремум по $\mu$ дают~1).

\textbf{2)} Пусть $0 < \lambda < \mu$. Из минимальности $x_\lambda$ и $x_\mu$:
\[
  \varphi(x_\lambda) + \lambda\alpha(x_\lambda)
    \le \varphi(x_\mu) + \lambda\alpha(x_\mu) , \qquad
  \varphi(x_\mu) + \mu\alpha(x_\mu)
    \le \varphi(x_\lambda) + \mu\alpha(x_\lambda) .
\]
Складывая: $(\mu-\lambda)\big(\alpha(x_\mu) - \alpha(x_\lambda)\big) \le 0$
$\Rightarrow$ $\alpha(x_\mu) \le \alpha(x_\lambda)$.

\textbf{3)} Из первого неравенства:
$\varphi(x_\lambda) - \varphi(x_\mu) \le
 \lambda\big(\alpha(x_\mu)-\alpha(x_\lambda)\big) \le 0$.

\textbf{4)}
$P(\lambda) \le \varphi(x_\mu) + \lambda\alpha(x_\mu)
 \le \varphi(x_\mu) + \mu\alpha(x_\mu) = P(\mu)$.

\textbf{5)} $P$ монотонна и ограничена сверху величиной $\min_S\varphi$ (п.~1)
$\Rightarrow$ предел равен супремуму. $\blacksquare$

\paragraph{Теорема о сходимости.}
\textit{Пусть решение $x_*$ существует, $\mu_k \to \infty$ и $\{x_{\mu_k}\}$
лежит в компакте. Тогда}
\[
  \min_{S}\varphi = \lim_k \varphi(x_{\mu_k}) , \qquad
  x_{\mu_k} \to x_* , \qquad
  \mu_k\,\alpha(x_{\mu_k}) \to 0 .
\]

\textsc{Док-во.} По лемме $\varphi(x_\mu) \le P(\mu) \le \varphi_*$, где
$\varphi_* = \min_S\varphi$, поэтому
\[
  \mu\,\alpha(x_\mu) = P(\mu) - \varphi(x_\mu)
   \le \varphi_* - \varphi(x_{\mu_1})
\]
ограничено сверху константой $\Rightarrow$ $\alpha(x_\mu) \to 0$ при
$\mu \to \infty$. Для предельной точки $\bar{x}$ по непрерывности
$\alpha(\bar{x}) = 0$ $\Rightarrow$ $\bar{x} \in S$ $\Rightarrow$
$\varphi(\bar{x}) \ge \varphi_*$. С другой стороны
$\varphi(\bar{x}) = \lim\varphi(x_{\mu_k}) \le \varphi_*$. Значит
$\varphi(\bar{x}) = \varphi_*$ — точка $\bar{x}$ оптимальна. Отсюда же
$P(\mu) \to \varphi_*$ и $\mu\alpha(x_\mu) = P(\mu) - \varphi(x_\mu) \to 0$.
$\blacksquare$

\paragraph{Алгоритм.}
\begin{csteps}
  \item Задать $\mu_1 > 0$, $\eps > 0$, $\lambda > 1$, начальную $x_1$ (может
        быть недопустимой), $k = 1$.
  \item Решить $\min_x\big[\varphi(x) + \mu_k\alpha(x)\big]$ любым методом
        безусловной минимизации (билеты~6--9); получить $x_{\mu_k}$.
  \item Если $\mu_k\alpha(x_{\mu_k}) < \eps$ — конец.
  \item Иначе $\mu_{k+1} = \lambda\mu_k$, $k := k+1$, к~шагу~2.
\end{csteps}

\textbf{Про критерий останова.} $\mu\alpha(x_\mu) < \eps$ означает
$\varphi_* - \varphi(x_\mu) \le \varphi_* - P(\mu) + \mu\alpha(x_\mu)$, где
$\varphi_* - P(\mu) \to 0$; практически это малое нарушение ограничений.

\paragraph{Внешние и внутренние штрафы.}
Здесь $\mu \to +\infty$ и точки могут быть недопустимыми. Во внутреннем
(барьерном, билет~16) методе наоборот: точки обязаны лежать строго внутри, а
параметр $\mu \to +0$.

\paragraph{Ограничения.}
При больших $\mu$ вспомогательная функция плохо обусловлена и минимизируется
тяжелее; её минимум лежит вне $S$; точное решение требует $\mu \to \infty$,
что недостижимо; слишком быстрый рост $\mu$ делает счёт неустойчивым,
слишком медленный — долгим.

\seefull{гл.~3, \S3 и \S11 (численный пример в \S11)}
